
en
posant pour tout @@M0@@
@@M1@@

\textbf{Remarques 5.2.2.} :
\begin{enumerate}[label=\alph*)]
\item Supposons que @@M2@@ soit le topos localement annelé des faisceaux zariskiens sur un schéma @@M3@@. Lorsque @@M4@@ est un diviseur de Cartier @@M5@@ sur @@M6@@ et
@@M7@@
alors on a
@@M8@@
où @@M9@@ est l'homomorphisme canonique de 1.2.3. Cela résulte immédiatement de la suite exacte
@@M10@@

\item Sous les hypothèses de 5.1 et 5.2 on a pour tout @@M11@@
@@M12@@
ce qui justifie la notation @@M13@@ introduite plus bas. En effet, on a
@@M14@@
\end{enumerate}

\subsection*{5.3.}
La composition de @@M15@@ avec l'épimorphisme canonique
@@M16@@
donne la «~première classe de Chern~» sur @@M17@@ :
@@M18@@

